\chapter{ALT (Alanine Aminotransferase)}

\section*{What Is This Test?}
Alanine aminotransferase (ALT), formerly called serum glutamic-pyruvic transaminase (SGPT), is an enzyme that catalyzes the reversible transfer of an amino group from alanine to alpha-ketoglutarate, producing pyruvate and glutamate. ALT is found in high concentrations in the cytoplasm of hepatocytes (liver cells), with much lower concentrations in kidney, heart, and skeletal muscle. Because ALT is predominantly located in the liver and is more specific to liver injury than AST, it is considered the most liver-specific aminotransferase. In hepatocellular injury, hepatocyte membrane damage causes leakage of ALT into the bloodstream. Serum ALT levels rise within hours of liver injury, peak at 1--3 days, and return to baseline over 1--3 weeks depending on the severity and persistence of injury. ALT has a serum half-life of approximately 47 hours (compared to approximately 17 hours for AST). The ALT concentration roughly parallels the degree of hepatocellular necrosis but does not distinguish between different causes of liver injury. ALT is measured as part of the comprehensive metabolic panel (CMP) and liver function tests (LFTs). The AST-to-ALT ratio (De Ritis ratio) aids in the differential diagnosis of liver disease: AST:ALT >2:1 suggests alcoholic liver disease (alcohol stimulates mitochondrial AST release and depletes pyridoxine, which is needed for ALT synthesis); AST:ALT <1 suggests viral hepatitis, NAFLD, or drug-induced liver injury.

\section*{Alternative Names}
SGPT; Serum Glutamic-Pyruvic Transaminase; Alanine Transaminase; GPT; Serum ALT; Liver Enzyme

\section*{Principle}
ALT is measured by a kinetic enzymatic spectrophotometric method based on the reaction: L-alanine + alpha-ketoglutarate $\rightarrow$ pyruvate + L-glutamate (catalyzed by ALT). The pyruvate produced is reduced to lactate by lactate dehydrogenase (LDH) with concurrent oxidation of NADH to NAD\textsuperscript{+}: Pyruvate + NADH + H\textsuperscript{+} $\rightarrow$ Lactate + NAD\textsuperscript{+}. The rate of decrease in absorbance at 340 nm (as NADH is consumed) is directly proportional to ALT activity in the sample. This method was standardized by the International Federation of Clinical Chemistry (IFCC) in 2002. The assay requires the addition of pyridoxal-5\textsuperscript{'}-phosphate (P5P, the active form of vitamin B6) to saturate ALT apoenzyme and ensure maximal activity (ALT requires P5P as a cofactor). Without P5P supplementation, ALT activity in pyridoxine-deficient patients (alcoholics, malnourished) may be falsely low. The test is performed at 37 degrees C on automated chemistry analyzers. Results are reported in IU/L (international units per liter). This test is NOT performed by hematology analyzers using impedance or flow cytometry methods.

\section*{History}
The transaminase enzymes (AST and ALT) were discovered by Arthur Karmen, Felix Wr\'{o}blewski, and John LaDue in the 1950s. In 1954, LaDue, Wr\'{o}blewski, and Karmen demonstrated that serum transaminase levels were elevated in patients with acute myocardial infarction (AST) and liver disease (AST and ALT). The clinical utility of the AST:ALT ratio (De Ritis ratio) was described by Fernando De Ritis in 1957. Wr\'{o}blewski and LaDue established the modern basis of serum enzyme diagnostics. The IFCC standardized method for ALT was published in 2002, and international standardization of enzyme assays has reduced inter-laboratory variability.

\section*{Why This Test Is Ordered}
\begin{itemize}
  \item Screening for hepatocellular liver injury as part of routine comprehensive metabolic panel
  \item Evaluation of patients with suspected hepatitis: acute viral hepatitis (HAV, HBV, HCV, HDV, HEV), autoimmune hepatitis, drug-induced liver injury, alcoholic hepatitis, ischemic hepatitis
  \item Monitoring of patients on potentially hepatotoxic medications: acetaminophen (acute overdose >150 mg/kg or >7.5 g in 24 hours), isoniazid, methotrexate, statins, valproic acid, amiodarone, ketoconazole, antiretroviral therapy, anti-TNF biologics, herbal supplements (green tea extract, kava, comfrey, chaparral)
  \item Screening for non-alcoholic fatty liver disease (NAFLD) and non-alcoholic steatohepatitis (NASH) in patients with obesity, metabolic syndrome, type 2 diabetes, and dyslipidemia
  \item Evaluation of unexplained fatigue, jaundice, right upper quadrant abdominal pain, hepatomegaly, or unexplained coagulopathy
  \item Preoperative assessment of liver function
  \item Monitoring of viral hepatitis treatment response and disease progression
\end{itemize}

\section*{Is This a Primary or Secondary Test}
ALT is a primary screening test as part of the CMP and is the most commonly used marker for hepatocellular injury. It is the cornerstone of liver function screening and is measured millions of times daily worldwide.

\section*{How This Test Is Performed}
A healthcare professional draws blood from a vein into a serum separator tube (gold-top) or lithium heparin tube (green-top). Approximately 3--5 mL of blood is collected. The sample is centrifuged and ALT activity is measured on an automated chemistry analyzer. Results are available within 30--60 minutes. Hemolysis must be avoided because red blood cells contain ALT at approximately 5--6 times the serum concentration, causing false elevation. The tube must be free of any visible hemolysis.

\section*{What Preparation Is Needed}
Fasting for 8--12 hours is recommended when ALT is ordered as part of a comprehensive metabolic panel (which includes glucose) or a lipid panel. ALT levels are affected by: (1) Recent alcohol consumption within 24--48 hours can elevate ALT. (2) Hepatotoxic medications: acetaminophen, isoniazid, methotrexate, valproic acid, amiodarone, statins (transaminase elevation in 1--3\% of patients, usually mild and self-limited), antiretrovirals. (3) Herbal supplements: green tea extract, kava, comfrey, chaparral, germander, jin bu huan. (4) Strenuous exercise within 24--48 hours can cause mild ALT elevation from muscle release (ALT is less concentrated in muscle than AST, so the AST rise is typically greater). (5) BMI and metabolic factors: obese and insulin-resistant patients have higher ALT at baseline. (6) Recent intramuscular injections can cause mild ALT release from muscle.

\section*{Contraindications and Precautions}
No absolute contraindications. Critical considerations: (1) Hemolysis is the most important pre-analytical interferent. Hemolyzed samples must be rejected for ALT measurement. (2) Serum or heparinized plasma are acceptable; EDTA or citrate plasma can chelate magnesium and calcium cofactors and reduce ALT activity. (3) Pyridoxine (vitamin B6) deficiency, common in alcoholics, malnourished patients, and the elderly, reduces ALT synthesis and activity because P5P is a required cofactor. In these patients, ALT may be falsely low despite significant liver injury. The IFCC method adds excess P5P to the assay to eliminate this effect, but not all laboratories use P5P-supplemented reagents. (4) Drugs falsely lowering ALT through assay interference are rare but include high concentrations of reducing agents (ascorbic acid). (5) ALT may be low in end-stage liver disease (``burned-out cirrhosis'') despite advanced fibrosis---correlate with synthetic function tests (INR, albumin, bilirubin). (6) A normal ALT does NOT exclude significant liver disease. Approximately 10--20\% of patients with chronic hepatitis C have persistently normal ALT; up to 79\% of patients with NAFLD may have ALT in the ``normal'' range (<40 IU/L). The upper limit of ``normal'' ALT has been lowered in some guidelines: 19--25 IU/L for females and 29--33 IU/L for males (Prati et al., 2002; Kwo et al., 2017 ACG guidelines).

\section*{Risks}
Minimal risk from venipuncture. Mild pain or bruising resolves in 1--2 days. Rare: hematoma, infection, nerve injury, vasovagal syncope.

\section*{What Happens After the Test}
Results available within 30--60 minutes. ALT interpretation requires concurrent AST, ALP, total bilirubin, and albumin. The pattern of liver test abnormalities guides the differential: (1) Predominantly ALT/AST elevation (hepatocellular pattern): viral hepatitis, NAFLD/NASH, autoimmune hepatitis, drug-induced liver injury, ischemic hepatitis, hemochromatosis, Wilson disease. (2) Predominantly ALP elevation (cholestatic pattern): bile duct obstruction, primary biliary cholangitis, primary sclerosing cholangitis, drug-induced cholestasis, infiltrative disease. (3) Mixed pattern: elevation of both aminotransferases and ALP. Severity grading: Mild (<5$\times$ ULN), Moderate (5--15$\times$ ULN), Severe (>15$\times$ ULN). ALT >1,000 IU/L is characteristic of acute viral hepatitis, acetaminophen toxicity, ischemic hepatitis (shock liver), and autoimmune hepatitis. ALT >3,000--5,000 IU/L: acetaminophen toxicity or ischemic hepatitis, which are the two most common causes of massive transaminase elevation. In acetaminophen toxicity, ALT typically rises to 3,000--10,000 IU/L at 24--72 hours post-ingestion and AST is similarly elevated.

\section*{Understanding Results}

\subsection*{Below Normal}
\begin{itemize}
  \item Low ALT is not a clinically significant finding and generally does not indicate liver disease. ALT below the lower reference limit is benign.
  \item ALT may be at the lower end of normal in elderly patients, patients with reduced muscle mass, and in patients with vitamin B6 (pyridoxine) deficiency.
  \item In end-stage liver disease (decompensated cirrhosis), ALT may fall into the normal range or become only mildly elevated despite advanced fibrosis because there are fewer viable hepatocytes to release ALT (``burned-out cirrhosis''). Always correlate with synthetic function markers (INR, albumin, bilirubin) and platelet count.
  \item Uremia inhibits ALT activity in vitro and can produce falsely low ALT in patients with CKD.
  \item A normal ALT does not exclude significant liver disease, including cirrhosis or chronic hepatitis C.
\end{itemize}

\subsection*{Above Normal}
\begin{itemize}
  \item \textbf{Mild elevation (1--3$\times$ ULN):} NAFLD (most common cause in developed countries---affects 25--35\% of the adult population; risk factors: obesity, type 2 diabetes, metabolic syndrome, hypertriglyceridemia). Chronic hepatitis C (approximately 80\% have mild-moderate ALT elevation; 10--20\% have persistently normal ALT). Alcoholic liver disease (AST:ALT >2:1 is classic). Medication-induced liver injury (statins, NSAIDs, antibiotics). Hemochromatosis (ALT may be mildly elevated before cirrhosis develops; ferritin >300 ng/mL males, >200 ng/mL females with transferrin saturation >45\%). Autoimmune hepatitis (ALT often in 200--1,000 IU/L range; anti-smooth muscle antibody, ANA, or anti-LKM1 positive; elevated IgG).
  \item \textbf{Moderate elevation (3--15$\times$ ULN):} Acute viral hepatitis (HAV, HBV, HCV, HDV, HEV). Autoimmune hepatitis. Wilson disease (especially in young patients <40 years; serum ceruloplasmin low, 24h urine copper elevated, Kayser-Fleischer rings on slit-lamp). Alpha-1 antitrypsin deficiency.
  \item \textbf{Marked elevation (>15$\times$ ULN, often >1,000 IU/L):} Acute viral hepatitis (HAV, HBV---ALT can reach 3,000--5,000 IU/L in the acute phase). Acetaminophen toxicity (the most common cause of acute liver failure in the US and UK; ALT often >3,000--10,000 IU/L with AST similarly elevated; acetaminophen level >150 mcg/mL at 4 hours post-ingestion is toxic). Ischemic hepatitis (shock liver, hypoperfusion): rapid rise and fall of ALT/AST with rapid normalization (<1 week) once perfusion is restored; ALT/AST >1,000 IU/L in setting of hypotension, sepsis, or cardiac arrest. Autoimmune hepatitis (ALT >1,000 IU/L in fulminant presentation). Drug-induced liver injury (isoniazid, valproic acid, halothane, herbal hepatotoxins).
  \item \textbf{The ALT/AST ratio (De Ritis ratio):} AST:ALT <1: viral hepatitis, NAFLD/NASH, drug-induced liver injury. AST:ALT >2:1: alcoholic liver disease (mitochondrial damage releases mitochondrial AST; also alcohol causes pyridoxine deficiency, which depletes ALT more than AST). AST:ALT >1 with cirrhosis: fibrosis reduces ALT production.
\end{itemize}

\section*{Normal Results}
\begin{itemize}
  \item \textbf{Adult males:} 10--40 IU/L (newer reference: <29--33 IU/L)
  \item \textbf{Adult females:} 7--35 IU/L (newer reference: <19--25 IU/L)
  \item \textbf{Children:} 5--45 IU/L (age and sex-dependent; slightly higher in adolescence)
  \item \textbf{Infants:} 10--55 IU/L (slightly higher than adults)
  \item \textbf{When to lower the upper limit of normal (Prati et al., 2002; ACG 2017):} Males >30 IU/L; Females >19 IU/L should prompt evaluation
  \item \textbf{Hepatocellular injury thresholds:} Mild: <5$\times$ ULN; Moderate: 5--15$\times$ ULN; Severe: >15$\times$ ULN; Massive: >1,000 IU/L
  \item \textbf{ALT half-life:} 47 hours
\end{itemize}

\section*{Follow-up Tests}
\begin{itemize}
  \item \textbf{AST, ALP, GGT, total bilirubin, albumin, INR/PT:} Complete liver panel for pattern recognition and synthetic function assessment.
  \item \textbf{AST:ALT ratio:} >2:1 suggests alcoholic liver disease; <1 suggests viral or NAFLD.
  \item \textbf{Viral hepatitis serologies:} HBsAg, anti-HBc IgM, anti-HAV IgM, anti-HCV, HCV RNA.
  \item \textbf{Abdominal ultrasound with Doppler:} First-line imaging for fatty liver, biliary obstruction, cirrhosis, portal hypertension. Screening for hepatocellular carcinoma in cirrhotic patients.
  \item \textbf{Serum ferritin, transferrin saturation, iron:} For hemochromatosis (ferritin >300 ng/mL males, >200 females; TSAT >45\%).
  \item \textbf{Serum ceruloplasmin, 24-hour urine copper:} For Wilson disease (ceruloplasmin <20 mg/dL, 24h urine Cu >100 mcg/day).
  \item \textbf{Autoimmune markers:} ANA, anti-smooth muscle antibody (ASMA), anti-LKM1, IgG for autoimmune hepatitis.
  \item \textbf{Alpha-1 antitrypsin level and phenotype:} For alpha-1 antitrypsin deficiency.
  \item \textbf{Acetaminophen level:} If drug-induced liver injury suspected; level at 4 hours post-ingestion determines N-acetylcysteine (NAC) therapy.
  \item \textbf{Liver biopsy:} Gold standard for staging fibrosis and diagnosing indeterminate liver disease.
  \item \textbf{FibroScan (transient elastography):} Non-invasive liver stiffness measurement; replaces biopsy for fibrosis staging in many scenarios.
\end{itemize}

\section*{References}
\begin{enumerate}
  \item Burtis CA, Ashwood ER, Bruns DE. Tietz Textbook of Clinical Chemistry and Molecular Diagnostics. 7th ed. Elsevier; 2023.
  \item Wallach J. Interpretation of Diagnostic Tests. 10th ed. Wolters Kluwer; 2022.
  \item Tierney LM, Saint S, Drazen JM. CURRENT Medical Diagnosis and Treatment. McGraw-Hill; 2023.
  \item Fischbach FT, Dunning MB. A Manual of Laboratory and Diagnostic Tests. 10th ed. Wolters Kluwer; 2023.
  \item Kwo PY, Cohen SM, Lim JK. ACG clinical guideline: evaluation of abnormal liver chemistries. Am J Gastroenterol. 2017;112(1):18--35.
  \item Prati D, Taioli E, Zanella A, et al. Updated definitions of healthy ranges for serum alanine aminotransferase levels. Ann Intern Med. 2002;137(1):1--10.
  \item De Ritis F, Coltorti M, Giusti G. An enzymic test for the diagnosis of viral hepatitis: the transaminase serum activities. Clin Chim Acta. 1957;2(1):70--74.
  \item Lee WM. Drug-induced hepatotoxicity. N Engl J Med. 2003;349(5):474--485.
\end{enumerate}

\clearpage
